%%%%%%%%%%%%%%%%%%%%%%%%%%%%%%%%%%%%%%%%%%%%%%%%%%%%%%%%%%%%%%%%%%%%%%%%%%%%%%%%
%
% Resume
% Author: Andrew Hard
% Date: August 19, 2016
%
%%%%%%%%%%%%%%%%%%%%%%%%%%%%%%%%%%%%%%%%%%%%%%%%%%%%%%%%%%%%%%%%%%%%%%%%%%%%%%%%

\documentclass{letter}

\usepackage{amssymb} % for black square
\usepackage[usenames, dvipsnames]{color}
%\definecolor{Maroon}{cmyk}{0, 0.87, 0.68, 0.32}
\definecolor{Maroon}{cmyk}{0, 0.9, 0.8, 0.5}
\usepackage{geometry} % for paper geometry
\usepackage{graphicx}
\usepackage{courier}
\usepackage{hyperref} % for hyperlinks
\hypersetup{hidelinks}
%\hypersetup{colorlinks=true,urlcolor=blue}
%\usepackage{helvetica}
\usepackage{newcent}
% Set the paper size and margin size (use a4paper, not legalpaper):
\geometry{a4paper, portrait, margin=0.5in}
\begin{document} 

\pagenumbering{gobble}

%%%%%%%%%%%%%%%%%%%%%%%%%%%%%%%%%%%%%%%%%%%%%%%%%%%%%%%%%%%%%%%%%%%%%%%%%%%%%%%%
\begin{tabular}{p{0.5\textwidth}p{0.5\textwidth}}
	\hfill \newline \href{https://ch.linkedin.com/in/andrew-hard-25b690a5}{\Huge{\color{Maroon}{Andrew Hard}}} \newline \LARGE{\textit{R\'{e}sum\'{e}}} \newline
	&
	\hfill \newline CERN PH Division, B32 RA-14, 1211 Gen\`{e}ve, Suisse \newline
	(+41) 076 30 88 007, (+1) 423 227 4106 \newline
	\href{mailto:ahard@cern.ch}{andrew.straiton.hard@cern.ch} \newline
	\href{https://github.com/rasumovsky}{\texttt{github.com/rasumovsky}}\\
\end{tabular}

%%%%%%%%%%%%%%%%%%%%%%%%%%%%%%%%%%%%%%%%%%%%%%%%%%%%%%%%%%%%%%%%%%%%%%%%%%%%%%%%
\begin{flushleft}
\Large{\textsc{\textbf{\color{Maroon}{Education}}}}
\vspace{1pt} % space between text and line
\hrule
\end{flushleft}

\begin{tabular}{p{0.15\textwidth}p{0.8\textwidth}}
	2016 (Expected)
	&
	\textbf{Doctor of Philosophy} in Physics \newline 
	\textit{University of Wisconsin}, Madison WI \newline
	Thesis: \textit{Searches and Discoveries with the $\gamma\gamma$ final state at ATLAS} \newline
	Advised by Prof. Sau Lan Wu \\
\\
	2010 
	& 
	\textbf{Bachelor of Arts} in Physics, Honors \newline 
	\textit{University of Chicago}, Chicago IL \newline
	Advised by Prof. Edward Blucher 
\\
\end{tabular}

%%%%%%%%%%%%%%%%%%%%%%%%%%%%%%%%%%%%%%%%%%%%%%%%%%%%%%%%%%%%%%%%%%%%%%%%%%%%%%%%
\begin{flushleft}
\Large{\textsc{\textbf{\color{Maroon}{Experience}}}}
\vspace{1pt} % space between text and line
\hrule
\end{flushleft}

%\renewcommand{\arraystretch}{0.5}% Wider
\begin{tabular}{p{0.15\textwidth}p{0.8\textwidth}}
	2011 - 2016
	&
	\textbf{Graduate Research Assistant}, \textit{Department of Physics, University of Wisconsin} \newline
	\vspace{-15pt}      % suppress blank line after tabbing
	\begin{itemize}
		\itemsep0em
%		\renewcommand{\labelitemi}{\scriptsize$\blacksquare$} 
		\renewcommand{\labelitemi}{\tiny$\blacksquare$}
    		\item Discovered Higgs boson, performed first measurements of couplings, spin, and mass 
		\item Contributed significantly to 19 papers \& notes, author on 250+ ATLAS publications
    		\item Organizer for first $h\rightarrow\gamma\gamma$ coupling result, statistical expert for multiple analyses
		\item Invented algorithm to spatially and temporally match CMOS chip hits at LBNL 
    		\item Developed analysis software with C++, ROOT, \& shell scripts
		\item Optimized physics searches in large phase spaces using massive $O(10^{9})$ datasets 
		\item Wrote and coordinated DoE funding reports for Wisconsin ATLAS Group
	\end{itemize}
\\
	2014 
	& 
	\textbf{Graduate Teaching Assistant}, \textit{Department of Physics, University of Wisconsin} \newline
	\vspace{-15pt}      % suppress blank line after tabbing
	\begin{itemize}
		\itemsep0em
		\renewcommand{\labelitemi}{\tiny$\blacksquare$} 
		\item Led discussions and labs on classical mechanics, electrodynamics, thermodynamics
    		\item Designed supplemental exercises and summary notes that boosted exam performances
	\end{itemize}
\\
	2010 - 2011
	&
	\textbf{CERN Technologist}, \textit{Enrico Fermi Institute, University of Chicago} \newline
	\vspace{-15pt}      % suppress blank line after tabbing
	\begin{itemize}
		\itemsep0em
		\renewcommand{\labelitemi}{\tiny$\blacksquare$} 
   		\item Electronic calibration expert for the ATLAS Experiment hadronic calorimeter
   		\item Developed \& maintained calibration software package using python and SQL
    		\item Documented, monitored and reported on detector status to collaboration
	\end{itemize}
\\
	2009 - 2010
	&
	\textbf{Undergraduate Research Assistant}, \textit{Enrico Fermi Institute, University of Chicago} \newline
	\vspace{-15pt}      % suppress blank line after tabbing
	\begin{itemize}
		\itemsep0em
		\renewcommand{\labelitemi}{\tiny$\blacksquare$} 
   		\item Developed particle detector simulation in C++, ROOT, and Geant4 
   		\item Constructed $\mu$ particle modules, worked in machine shop, tested electronics 
	\end{itemize}
\end{tabular}

\vspace{-10pt}

%%%%%%%%%%%%%%%%%%%%%%%%%%%%%%%%%%%%%%%%%%%%%%%%%%%%%%%%%%%%%%%%%%%%%%%%%%%%%%%%
\begin{flushleft}
\Large{\textsc{\textbf{\color{Maroon}{Skills}}}}
\vspace{1pt} % space between text and line
\hrule
\end{flushleft}

\begin{tabular}{p{0.15\textwidth}p{0.8\textwidth}}
	{\bf Scientific} 
	&
	Physics, Statistics, Simulation, High Throughput Computing, Numerical Methods, Data Structures, Machine Learning, Collaborative Research, Public Presentation \newline
\\
	{\bf Programming} 
	&
	C++, Python, Java, \LaTeX{}, Unix/Linux shell scripting, ROOT, Matlab, SQL, TensorFlow \newline
\\ 
	{\bf Languages}
	&
	English (native), French (basic oral and written communication)
\end{tabular}

%%%%%%%%%%%%%%%%%%%%%%%%%%%%%%%%%%%%%%%%%%%%%%%%%%%%%%%%%%%%%%%%%%%%%%%%%%%%%%%%
\begin{flushleft}
\Large{\textsc{\textbf{\color{Maroon}{Volunteering \& Outreach}}}}
\vspace{1pt} % space between text and line
\hrule
\end{flushleft}
	\vspace{-10pt}      % suppress blank line after tabbing
\begin{tabular}{p{0.15\textwidth}p{0.8\textwidth}}
	\href{https://twiki.cern.ch/twiki/pub/AtlasPublic/HiggsPublicResults//Hgg-FloatingScale-Short2.gif}{\raisebox{-0.95\height}{\includegraphics[scale=0.16]{AnimationQR.jpg}}}
	&
	\begin{itemize}
		\itemsep0em
		\renewcommand{\labelitemi}{\tiny$\blacksquare$} 
		\item Newtonian physics demonstration for Chicago Public Library \hfill 2016
		\item US voter outreach \& registration at CERN \hfill 2016
		\item Discussed research funding with U.S. lawmakers in Washington D.C. \hfill 2014 \& 2015
		\item Created GIF visualizations of Higgs boson discovery data \hfill 2013
		\item Visited classrooms at the Chattanooga School for the Arts \& Sciences \hfill 2012
	\end{itemize}	
\end{tabular}

\vspace{-10pt}

%%%%%%%%%%%%%%%%%%%%%%%%%%%%%%%%%%%%%%%%%%%%%%%%%%%%%%%%%%%%%%%%%%%%%%%%%%%%%%%%
\begin{flushleft}
\Large{\textsc{\textbf{\color{Maroon}{Awards}}}}
\vspace{1pt} % space between text and line
\hrule
\end{flushleft}

\begin{tabular}{p{0.15\textwidth}p{0.8\textwidth}}
	2015
	&
	\textbf{Teaching Assistant Rookie of the Year}, \textit{Department of Physics, University of Wisconsin} 
\\
	2013 \& 2014
	& 
	\textbf{Lighting Round Winner}, \textit{US LHC User's Association Annual Meeting}
\\
\end{tabular}

\end{document}